% !TEX program = xelatex
\documentclass[12pt,a4paper]{ctexart}

% ==================== 字体设置 ====================
% 中文主字体：宋体（启用伪粗体解决加粗问题）
\setCJKmainfont{simsun.ttc}[
    Path        = C:/Windows/fonts/,
    AutoFakeBold = 2,
    ItalicFont   = simkai.ttf,
]
% 中文无衬线：黑体
\setCJKsansfont{simhei.ttf}[
    Path        = C:/Windows/fonts/,
    AutoFakeBold = 2,
]
% 中文等宽：仿宋
\setCJKmonofont{simfang.ttf}[
    Path        = C:/Windows/fonts/,
]

% 英文主字体
\setmainfont{Latin Modern Roman}
\setsansfont{Latin Modern Sans}
\setmonofont{Latin Modern Mono}

% ==================== 页面布局 ====================
\usepackage[top=2.5cm, bottom=2.5cm, left=2.8cm, right=2.8cm,
            headheight=14pt]{geometry}

% ==================== 包 ====================
\usepackage{hyperref}
\usepackage{graphicx}
\usepackage{booktabs}
\usepackage{longtable}
\usepackage{enumitem}
\usepackage{fancyhdr}
\usepackage{xcolor}
\usepackage{titlesec}

% ---------- listings (代码高亮) ----------
\usepackage{listings}
\lstset{
    basicstyle       = \small\ttfamily,
    backgroundcolor  = \color{gray!5},
    frame            = single,
    rulecolor        = \color{gray!40},
    breaklines       = true,
    numbers          = left,
    numbersep        = 6pt,
    numberstyle      = \tiny\color{gray},
    tabsize          = 4,
    showstringspaces = false,
    extendedchars    = true,
    literate         = {→}{{$\rightarrow$}}1
                       {×}{{$\times$}}1,
}

% ==================== 颜色 ====================
\definecolor{Accent}{RGB}{0,102,204}

% ==================== 页眉页脚 ====================
\pagestyle{fancy}
\fancyhf{}
\fancyhead[L]{\small\sffamily EdgeHealth-ML}
\fancyhead[R]{\small\sffamily\thepage}
\renewcommand{\headrulewidth}{0.3pt}
\setlength{\headheight}{14pt}

% ==================== 标题格式 ====================
\titleformat{\section}
    {\Large\bfseries\color{Accent}}
    {}{0em}{}[\vspace{0.1em}{\color{Accent}\hrule}]
\titleformat{\subsection}
    {\large\bfseries}
    {}{0em}{}
\titleformat{\subsubsection}
    {\normalsize\bfseries}
    {}{0em}{}

% ==================== 间距 ====================
\setlength{\parskip}{0.8ex plus 0.3ex minus 0.2ex}
\setlength{\parindent}{0pt}

% ==================== 命令 ====================
\newcommand{\file}[1]{\texttt{#1}}
\newcommand{\code}[1]{\texttt{#1}}

% =================================================================
\begin{document}

% -------------------- 封面 --------------------
\begin{center}
    {\huge\bfseries EdgeHealth-ML}

    \vspace{0.3em}
    {\LARGE 端侧多模态人体活动识别系统}

    \vspace{0.5em}
    {\large 完整项目分析与代码详解}

    \vspace{2em}
    {\normalsize 吴金骏 \textbar{} WU, Jinjun}

    \vspace{0.2em}
    {\normalsize Hong Kong University of Science and Technology}

    \vspace{0.2em}
    {\normalsize \today}
\end{center}

% -------------------- 摘要 --------------------
\begin{center}
\fbox{\begin{minipage}{0.92\textwidth}
\small
\textbf{摘要：} 本文档对 EdgeHealth-ML 项目进行完整技术分析，涵盖项目背景、模型架构（Vision+IMU 双分支融合网络）、数据处理流水线（STFT 时频谱图生成）、训练流程（30 epochs，最佳准确率 87.61\%）、INT8 量化与 ONNX 部署（ONNX Runtime 推理延迟 4.83ms），以及全部 6 个 Python 源文件的逐行代码详解。项目面向移动/边缘设备的实时人体活动识别，对标 HKUST MINSys Lab 研究方向。
\end{minipage}}
\end{center}

% =================================================================
\newpage
\section{项目概述}

\subsection{背景与目标}

EdgeHealth-ML 是一个面向\textbf{移动/边缘设备}的\textbf{多模态人体活动识别}（Human Activity Recognition, HAR）系统。利用智能手机自带的\textbf{加速度计 + 陀螺仪}（IMU 传感器），在设备端实时识别用户的日常活动（走路、上楼、下楼、坐着、站着、躺着），无需联网。项目对标 HKUST MINSys Lab（欧阳晓敏教授）的研究方向。

\subsection{技术亮点}

\begin{enumerate}[leftmargin=*]
    \item \textbf{双分支异构融合：} Vision（2D CNN 处理频谱图）+ IMU（1D CNN 处理时序），互补特征通过拼接融合
    \item \textbf{端到端部署：} PyTorch 训练 $\rightarrow$ INT8 量化 $\rightarrow$ ONNX 导出 $\rightarrow$ ONNX Runtime 推理
    \item \textbf{轻量化设计：} MobileNetV3-Small 作为视觉主干，总参数量仅 120 万
    \item \textbf{频谱图转换：} 将 1D 时序信号通过 STFT 转为 2D 时频谱图，复用成熟图像分类架构
    \item \textbf{量化无损：} INT8 动态量化后准确率仅下降 0.03\%，ONNX Runtime 推理加速 11 倍
\end{enumerate}

\subsection{数据集}

UCI HAR Dataset：
\begin{itemize}
    \item 30 名受试者，佩戴智能手机执行 6 种活动
    \item 传感器采样率 50Hz，每个样本 128 个时间步（2.56 秒窗口）
    \item 训练集 7,352 样本，测试集 2,947 样本
    \item 6 类活动：walking, walking\_upstairs, walking\_downstairs, sitting, standing, laying
\end{itemize}

% =================================================================
\section{模型架构}

\subsection{整体结构}

模型由 Vision 分支（MobileNetV3-Small）处理时频谱图、IMU 分支（3 层 1D-CNN）处理原始传感器信号、融合 MLP 分类三部分组成。

\begin{verbatim}
┌─────────────────────────────┐    ┌─────────────────────────────┐
│  Vision 分支 (VisionEncoder) │    │  IMU 分支 (IMUEncoder)       │
│  MobileNetV3-Small           │    │  3 层 1D-CNN                │
│  输入: (B, 3, 224, 224)      │    │  输入: (B, 6, 128)           │
│  输出: (B, 576)              │    │  输出: (B, 128)              │
└──────────────┬──────────────┘    └──────────────┬──────────────┘
               │                                  │
               └────────────────┬─────────────────┘
                                ▼
                   ┌───────────────────────┐
                   │  Fusion MLP            │
                   │  concat(576+128=704)   │
                   │  → Linear(704→256)     │
                   │  → BN + ReLU + Dropout │
                   │  → Linear(256→6)       │
                   └───────────┬───────────┘
                               ▼
                   ┌───────────────────────┐
                   │  6 类活动分类输出       │
                   └───────────────────────┘
\end{verbatim}

\subsection{VisionEncoder：频谱图编码器}

\begin{itemize}
    \item \textbf{主干网络：} MobileNetV3-Small，使用 ImageNet 预训练权重（迁移学习）
    \item \textbf{输入：} 3 通道时频谱图 $(B, 3, 224, 224)$
    \item \textbf{处理流程：} \code{backbone.features}（去掉分类头）$\rightarrow$ \code{AdaptiveAvgPool2d(1)}（全局平均池化）
    \item \textbf{输出：} 576 维特征向量 $(B, 576)$
\end{itemize}

MobileNetV3-Small 使用深度可分离卷积（depthwise separable convolution）和 SE（Squeeze-and-Excitation）注意力模块，在保持较高精度的同时大幅减少计算量。

\subsection{IMUEncoder：时序信号编码器}

\begin{itemize}
    \item \textbf{结构：} 3 层 1D 卷积 + BatchNorm + ReLU
    \item \textbf{输入：} 原始 6 通道传感器信号 $(B, 6, 128)$——3 轴加速度 + 3 轴陀螺仪，128 个时间步
    \item \textbf{维度变化：}
    \begin{verbatim}
    (B, 6, 128)   → Conv1d(6→64, k=5, s=2)  → (B, 64, 64)
    (B, 64, 64)   → Conv1d(64→128, k=5, s=2) → (B, 128, 32)
    (B, 128, 32)  → Conv1d(128→128, k=5, s=2) → (B, 128, 16)
    (B, 128, 16)  → AdaptiveAvgPool1d(1)      → (B, 128)
    \end{verbatim}
    \item \textbf{卷积核大小 5：} 捕捉 5 个时间步（0.1 秒 @ 50Hz）的局部运动模式
    \item \textbf{步长 2：} 逐步降采样，增大感受野
\end{itemize}

\subsection{融合策略：晚期融合}

\begin{itemize}
    \item \textbf{融合方式：} Concat 拼接 $\rightarrow$ 704 维 $(576 + 128)$
    \item \textbf{MLP 结构：} \code{Linear(704→256) → BatchNorm → ReLU → Dropout(0.3) → Linear(256→6)}
    \item \textbf{Dropout(0.3)：} 唯一的正则化手段，防止过拟合
\end{itemize}

\textbf{为什么选择晚期融合？} 两个模态的物理意义完全不同——频谱图是频域二维图像，IMU 是时域一维信号。先让各自分支独立提取模态特定特征，再在高维语义层融合，避免模态间干扰。

% =================================================================
\section{数据处理流水线}

\subsection{原始数据格式}

UCI HAR Dataset 的原始格式是空格分隔的纯文本文件：

\begin{itemize}
    \item \file{body\_acc\_x\_train.txt}：身体加速度 X 轴，形状 $(7352, 128)$
    \item \file{y\_train.txt}：训练标签，1--6 映射到 0--5
\end{itemize}

共 9 个信号文件（6 个身体传感器 + 3 个总加速度），加 2 个标签文件。

\subsection{\file{\_load\_txt()} — 文本读取}

\begin{lstlisting}[language=Python]
def _load_txt(path):
    with open(path, "r") as f:
        return np.array(
            [[float(x) for x in line.strip().split()] for line in f]
        )
\end{lstlisting}

逐行读取 → 去除首尾空白 → 按空格分割 → 转为浮点数 → 返回 $(N, 128)$ 的 numpy 数组。

\subsection{双模态数据生成}

每个原始样本的传感器数据被处理成两种互补的输入形式：

\subsubsection{模态 1：时频谱图（Vision 分支输入）}

\begin{center}
加速度信号 $(3, 128) \rightarrow$ STFT $\rightarrow$ |幅度| $\rightarrow$ log(1+x) $\rightarrow$ resize(224,224) $\rightarrow$ $(3, 224, 224)$
\end{center}

STFT 参数：\code{nperseg=64}（1.28 秒窗口）、\code{noverlap=48}（75\% 重叠）、\code{window="hann"}（汉宁窗）。

关键设计：
\begin{itemize}
    \item 3 轴加速度各自独立做 STFT → 3 张谱图沿 channel 维度堆叠
    \item \code{log1p} 压缩动态范围：原始幅度 $\approx [0, 100]$，取 log 后 $\approx [0, 5]$，利于神经网络学习
    \item \code{scipy.ndimage.zoom}（order=1，双线性插值）缩放到 224×224
    \item 最后 \code{Normalize(mean=0.5, std=0.5)} 映射到 $[-1, 1]$
\end{itemize}

\subsubsection{模态 2：原始 IMU 信号（IMU 分支输入）}

\begin{itemize}
    \item 3 轴加速度 + 3 轴陀螺仪 = 6 通道 × 128 时间步 → 直接作为 tensor
    \item \textbf{不做额外处理}，保留原始时序信息
    \item 加速度反映姿态和运动方向，陀螺仪反映旋转角速度
\end{itemize}

\subsection{HARDataLoader 类}

继承 \code{torch.utils.data.Dataset}，\code{\_\_getitem\_\_} 返回三元组 \code{(spectrogram, imu\_signal, label)}。

\begin{lstlisting}[language=Python]
class HARDataLoader(Dataset):
    def __getitem__(self, idx):
        # 模态 1：时频谱图
        acc_signals = np.stack([body_acc_x[idx], ...], axis=0)  # (3, 128)
        spec = _build_spectrogram_image(acc_signals, 224)       # (3, 224, 224)
        spec = torch.from_numpy(spec)
        spec = spec_norm(spec)  # Normalize

        # 模态 2：原始 IMU 信号
        imu = np.stack([body_acc_x[idx], ..., body_gyro_z[idx]], axis=0)
        imu = torch.from_numpy(imu).float()  # (6, 128)

        return spec, imu, int(labels[idx])
\end{lstlisting}

\subsection{create\_dataloaders() 工厂函数}

\begin{lstlisting}[language=Python]
def create_dataloaders(data_dir, batch_size=32):
    train_loader = DataLoader(train_ds, batch_size=batch_size, shuffle=True)
    test_loader  = DataLoader(test_ds,  batch_size=batch_size, shuffle=False)
    return train_loader, test_loader, len(LABELS)
\end{lstlisting}

\code{shuffle=True} 训练时打乱，\code{False} 测试时保持顺序；\code{num\_workers=0} 单进程（Windows 兼容）。

% =================================================================
\section{训练流程}

\subsection{超参数配置}

\begin{table}[h]
\centering
\begin{tabular}{@{}lll@{}}
\toprule
\textbf{参数} & \textbf{值} & \textbf{说明} \\
\midrule
优化器          & AdamW             & weight\_decay=$10^{-4}$ \\
初始学习率       & 0.001            & AdamW 初始步长 \\
学习率调度       & CosineAnnealingLR & 余弦衰减到 0 \\
损失函数         & CrossEntropyLoss  & 多分类交叉熵 \\
Batch Size      & 32               & 每批 32 个样本 \\
Epochs          & 30               & 完整遍历 30 次 \\
设备            & CPU              & Intel Core（无 GPU）\\
\bottomrule
\end{tabular}
\end{table}

\subsection{余弦退火学习率}

学习率按余弦曲线从初始值平滑衰减到 0：

\[
\mathrm{lr}(t) = \frac{\mathrm{lr}_0}{2}\left(1 + \cos\frac{t\pi}{T}\right)
\]

$t$ 为当前 epoch，$T=30$ 为总 epoch 数。后期极小学习率帮助模型收敛到更优局部最小值。

\subsection{训练循环}

每个 epoch 执行：
\begin{enumerate}
    \item \code{model.train()}：启用 Dropout 和 BatchNorm 训练模式
    \item \code{train\_one\_epoch()}：遍历训练集，前向 → 计算损失 → 反向传播 → 更新权重
    \item \code{evaluate()}：测试集评估（\code{@torch.no\_grad()} 禁用梯度计算）
    \item \code{scheduler.step()}：衰减学习率
    \item 若 test\_acc > best\_acc，保存 \file{best\_model.pth}
\end{enumerate}

\subsection{训练日志（完整 30 epochs）}

\begin{longtable}{@{}cccc@{}}
\toprule
\textbf{Epoch} & \textbf{Train Loss} & \textbf{Train Acc} & \textbf{Test Acc} \\
\midrule
\endhead
\midrule
\multicolumn{4}{r}{\small（续下页）}
\endfoot
\bottomrule
\endlastfoot
 1 & 0.5088 & 0.775 & 0.251 * \\
 2 & 0.3534 & 0.847 & 0.707 * \\
 3 & 0.3119 & 0.865 & 0.713 * \\
 4 & 0.2753 & 0.889 & 0.687   \\
 5 & 0.2496 & 0.894 & 0.649   \\
 6 & 0.2205 & 0.910 & 0.717 * \\
 7 & 0.1979 & 0.918 & 0.692   \\
 8 & 0.1950 & 0.918 & 0.716   \\
 9 & 0.1621 & 0.936 & 0.795 * \\
10 & 0.1293 & 0.948 & 0.820 * \\
11 & 0.1083 & 0.959 & 0.730   \\
12 & 0.0947 & 0.963 & 0.832 * \\
13 & 0.0731 & 0.974 & 0.799   \\
14 & 0.0524 & 0.980 & 0.785   \\
15 & 0.0394 & 0.986 & 0.841 * \\
16 & 0.0318 & 0.988 & 0.795   \\
17 & 0.0270 & 0.991 & 0.806   \\
18 & 0.0138 & 0.995 & 0.861 * \\
19 & 0.0116 & 0.997 & 0.842   \\
20 & 0.0133 & 0.995 & 0.835   \\
21 & 0.0080 & 0.998 & 0.800   \\
22 & 0.0030 & 1.000 & 0.854   \\
23 & 0.0040 & 0.999 & 0.827   \\
24 & 0.0054 & 0.999 & 0.847   \\
25 & 0.0024 & 0.999 & \textbf{0.876} * \\
26 & 0.0021 & 1.000 & 0.871   \\
27 & 0.0013 & 1.000 & 0.871   \\
28 & 0.0017 & 1.000 & 0.874   \\
29 & 0.0012 & 1.000 & 0.873   \\
30 & 0.0017 & 1.000 & 0.873   \\
\caption{训练日志（* 表示打破历史最佳测试准确率）}
\end{longtable}

\subsection{最终评估结果}

最佳模型在 Epoch 25，测试准确率 \textbf{87.61\%}。

\begin{table}[h]
\centering
\begin{tabular}{@{}lcccc@{}}
\toprule
\textbf{活动类别} & \textbf{Precision} & \textbf{Recall} & \textbf{F1} & \textbf{样本数} \\
\midrule
walking            & 0.9841 & 1.0000 & 0.9920 & 496 \\
walking\_upstairs   & 0.9638 & 0.9618 & 0.9628 & 471 \\
walking\_downstairs & 0.9758 & 0.9595 & 0.9676 & 420 \\
sitting            & 0.7140 & 0.6456 & 0.6781 & 491 \\
standing           & 0.7522 & 0.8158 & 0.7827 & 532 \\
laying             & 0.8887 & 0.8920 & 0.8903 & 537 \\
\midrule
\textbf{总体}      & \multicolumn{3}{c}{\textbf{0.8761}} & 2947 \\
\bottomrule
\end{tabular}
\caption{分类报告（测试集）}
\end{table}

\textbf{结论：}
\begin{itemize}
    \item \textbf{动态活动（walking 系列）}识别近乎完美——walking 100\% 召回率，上下楼 >96\%
    \item \textbf{静态活动（sitting/standing）}存在混淆——两者传感器模式极为相似，这是 HAR 经典难题
    \item \textbf{laying} 识别较好（89\%）——躺下时加速度方向与直立明显不同
\end{itemize}

\subsection{混淆矩阵}

\begin{verbatim}
              walking  up  down  sit  stand  lay
walking          496    0     0    0      0    0
walking_up         8  453    10    0      0    0
walking_down       0   17   403    0      0    0
sitting            0    0     0  317    134   40
standing           0    0     0   78    434   20
laying             0    0     0   49      9  479
\end{verbatim}

主要混淆：sitting $\leftrightarrow$ standing（212 样本）、sitting $\rightarrow$ laying（40 样本）。

% =================================================================
\section{模型优化：量化与 ONNX 部署}

\subsection{FP32 ONNX 导出}

训练好的 PyTorch 模型（\file{.pth}）只能被 PyTorch 加载。ONNX 导出使其可在任何支持 ONNX Runtime 的设备运行（手机、嵌入式等）。

\textbf{导出原理：}
\begin{enumerate}
    \item 用随机数据（dummy input）运行模型，\code{torch.export} 追踪所有算子调用
    \item PyTorch 计算图转为 ONNX 中间表示（IR）
    \item 序列化为 \file{.onnx} 文件
\end{enumerate}

\textbf{关键参数：}
\begin{itemize}
    \item \texttt{dynamic\_axes=\{0: "batch"\}}：batch 维度动态化，推理时可传入任意 batch size
    \item \code{opset\_version=18}：匹配 PyTorch 2.12 的 ONNX 算子集
\end{itemize}

\subsection{INT8 动态量化}

\textbf{原理：}
\begin{verbatim}
原始:  权重(FP32) × 输入(FP32) = 输出(FP32)
量化:  [权重(INT8) → 反量化(FP32)] × 输入(FP32) = 输出(FP32)
\end{verbatim}

\begin{itemize}
    \item 仅量化 \code{nn.Linear} 和 \code{nn.Conv1d} 的权重——激活值保持 FP32
    \item 计算时临时反量化，精度损失极小
    \item 权重量化减少内存占用约 25\%
\end{itemize}

\subsection{量化效果对比}

\begin{table}[h]
\centering
\begin{tabular}{@{}lccc@{}}
\toprule
\textbf{指标} & \textbf{FP32} & \textbf{INT8} & \textbf{变化} \\
\midrule
准确率      & 87.61\%      & 87.58\%       & $-$0.03\% \\
推理延迟     & 53.91 ms     & 52.66 ms      & 1.02× \\
模型大小     & 4,972 KB     & 4,441 KB      & 1.12× \\
参数量       & 1,235,302    & 1,053,280     & $-$14.7\% \\
\bottomrule
\end{tabular}
\caption{FP32 vs INT8 量化对比}
\end{table}

\subsection{ONNX Runtime 推理性能}

\begin{table}[h]
\centering
\begin{tabular}{@{}lc@{}}
\toprule
\textbf{指标} & \textbf{结果} \\
\midrule
平均延迟   & 4.83 ms \\
标准差     & 2.53 ms \\
最小延迟   & 2.73 ms \\
最大延迟   & 21.80 ms \\
P95 延迟   & 9.44 ms \\
吞吐量     & 207.0 次/秒 \\
\bottomrule
\end{tabular}
\caption{ONNX Runtime 推理基准（FP32 ONNX，100 次迭代）}
\end{table}

\textbf{ONNX Runtime 比 PyTorch 原生推理快 11 倍}（4.83ms vs 53.91ms），原因：
\begin{enumerate}
    \item ONNX Runtime 对计算图做了算子融合和常量折叠优化
    \item PyTorch eager 模式有 Python 解释器开销，ONNX Runtime 纯 C++ 执行
    \item 批量推理时 ONNX Runtime 多线程并行效率更高
\end{enumerate}

% =================================================================
\section{项目文件结构与代码详解}

\subsection{目录结构}

\begin{verbatim}
EdgeHealth-ML/
├── data/
│   ├── download_har.py          # 多镜像下载 UCI HAR 数据集
│   └── UCI_HAR_Dataset/         # 原始传感器数据（txt 格式）
├── models/
│   └── multimodal_model.py      # VisionEncoder + IMUEncoder + Fusion
├── utils/
│   └── data_loader.py           # STFT 谱图生成 + 数据加载
├── train.py                     # 训练主脚本
├── quantize_export.py           # 量化 + ONNX 导出
├── inference_benchmark.py       # ONNX Runtime 推理基准测试
├── checkpoints/                 # 模型权重 + ONNX 文件
├── runs/                        # TensorBoard 训练日志
└── requirements.txt             # 依赖清单
\end{verbatim}

\subsection{执行顺序}

\begin{center}
\file{download\_har.py} $\rightarrow$ \file{data\_loader.py} $\rightarrow$ \file{multimodal\_model.py} \\
$\downarrow$ \\
\file{inference\_benchmark.py} $\leftarrow$ \file{quantize\_export.py} $\leftarrow$ \file{train.py}
\end{center}

\subsection{\file{data/download\_har.py} —— 数据集下载}

\textbf{功能：} 下载并解压 UCI HAR Dataset。被 \file{train.py} 和 \file{quantize\_export.py} 调用。

\textbf{关键设计：}
\begin{itemize}
    \item \textbf{多镜像容错}（\code{HAR\_URLS} 列表）：优先 UCI 官方源，失败则切换 GitHub 镜像
    \item \textbf{幂等性}：如数据集已解压，直接返回路径（\code{if os.path.exists(EXTRACT\_DIR)}）
    \item \textbf{完整性验证}：下载后用 \code{zipfile.ZipFile} 打开验证合法性，损坏则删除重试
    \item \textbf{清理}：解压后删除 zip 文件释放空间
\end{itemize}

\textbf{核心函数 \code{download\_and\_extract()}：}

\begin{lstlisting}[language=Python]
def download_and_extract():
    if os.path.exists(EXTRACT_DIR):        # 幂等性检查
        return EXTRACT_DIR
    for url in HAR_URLS:                   # 遍历镜像
        try:
            urlretrieve(url, zip_path)     # 下载
            with zipfile.ZipFile(zip_path) as test: pass  # 验证
            break
        except Exception:
            if os.path.exists(zip_path): os.remove(zip_path)
            continue                       # 下一个镜像
    else:
        sys.exit(1)                        # 全部失败 → 退出
    with zipfile.ZipFile(zip_path) as zf:
        zf.extractall(DATA_DIR)            # 解压
    os.remove(zip_path)                    # 清理
    return EXTRACT_DIR
\end{lstlisting}

\subsection{\file{utils/data\_loader.py} —— 数据加载与预处理}

\textbf{功能：} 从原始 txt 文件读取传感器数据，转化为双模态输入。

\textbf{关键函数：}
\begin{itemize}
    \item \code{\_load\_txt(path)}：读取空格分隔文本 → numpy 数组
    \item \code{\_signal\_to\_spectrogram(signal\_1d, nperseg=64, noverlap=48)}：STFT → |幅度| → log1p → 2D 谱图
    \item \code{\_build\_spectrogram\_image(signals\_3ch, 224)}：3 通道各做 STFT → stack → zoom 缩放 → normalize
    \item \code{\_resize\_2d(arr, h, w)}：\code{scipy.ndimage.zoom} 双线性插值
\end{itemize}

\textbf{核心类 \code{HARDataLoader(Dataset)}：} \code{\_\_getitem\_\_} 返回 \code{(spectrogram, imu\_signal, label)}。

\subsection{\file{models/multimodal\_model.py} —— 模型架构}

\textbf{功能：} 定义双分支多模态融合网络的全部组件。

\textbf{三个核心类：}

\textbf{1. \code{IMUEncoder}（1D-CNN 时序编码器）—— 行 27--55}
\begin{itemize}
    \item 3 层 Conv1d，每层 kernel=5, stride=2, padding=2
    \item 维度流：$(B,6,128) \rightarrow (B,64,64) \rightarrow (B,128,32) \rightarrow (B,128,16) \rightarrow (B,128)$
    \item \code{AdaptiveAvgPool1d(1)} 全局平均池化压缩时间维度
\end{itemize}

\textbf{2. \code{VisionEncoder}（频谱图编码器）—— 行 58--76}
\begin{itemize}
    \item \code{mobilenet.mobilenet\_v3\_small(weights="DEFAULT")} 加载预训练
    \item \code{backbone.features} 去掉分类头，\code{AdaptiveAvgPool2d(1)} 全局池化
    \item 输出 576 维——MobileNetV3-Small 的标准特征维度
\end{itemize}

\textbf{3. \code{MultimodalHARModel}（融合网络）—— 行 79--110}
\begin{itemize}
    \item 包含 vision\_encoder + imu\_encoder + fusion MLP
    \item Fusion: Linear(704→256) → BN → ReLU → Dropout(0.3) → Linear(256→6)
    \item \code{forward()}：分别编码 → \code{torch.cat} 拼接 → 分类
\end{itemize}

\textbf{4. \code{count\_parameters()} —— 行 116--119}
\begin{itemize}
    \item 统计总参数和可训练参数（本项目 1,235,302 个全部可训练）
\end{itemize}

\subsection{\file{train.py} —— 训练脚本}

\textbf{功能：} 编排完整训练流程。

\textbf{命令行参数（\code{parse\_args()}）：}
\code{--epochs 30}、\code{--batch\_size 32}、\code{--lr 0.001}、\code{--weight\_decay 1e-4}、\code{--device auto}

\textbf{关键函数：}
\begin{itemize}
    \item \code{train\_one\_epoch()}：前向 → 损失 → 反向 → 更新，返回平均 loss 和 accuracy
    \item \code{evaluate()}：\code{@torch.no\_grad()} 禁用梯度，返回 loss + acc + 所有预测和标签
    \item \code{main()}：4 步流程——
    \begin{enumerate}
        \item 准备数据（下载 + 创建 DataLoader）
        \item 构建模型（实例化 + 打印参数量）
        \item 训练循环（30 epochs + CosineAnnealingLR + TensorBoard + 保存最佳模型）
        \item 最终评估（\code{classification\_report} + \code{confusion\_matrix}）
    \end{enumerate}
\end{itemize}

\textbf{优化器选择：AdamW} —— 与标准 Adam 的区别在于 weight\_decay 直接作用于权重而非通过梯度，提供更好的正则化效果。

\subsection{\file{quantize\_export.py} —— 量化与 ONNX 导出}

\textbf{功能：} 将训练好的 PyTorch 模型压缩并导出为跨平台 ONNX 格式。

\textbf{关键函数：}
\begin{itemize}
    \item \code{export\_onnx(model, filepath)}：\code{torch.onnx.export()} 导出，支持动态 batch
    \item \code{quantize\_model(model)}：\code{torch.quantization.quantize\_dynamic()} INT8 动态量化
    \item \code{compare\_models(fp32, int8, loader)}：准确率 + 延迟 + 大小三方对比
    \item \code{main()}：4 步——加载数据 → 加载模型 → 导出 FP32 ONNX → INT8 量化 + 对比
\end{itemize}

\textbf{兼容性说明：} PyTorch 2.12+ 基于 \code{torch.export} 的 ONNX 导出器不支持量化模型。FP32 ONNX 导出成功（397 KB），INT8 ONNX 通过 try/except 跳过。

\subsection{\file{inference\_benchmark.py} —— 推理性能测试}

\textbf{功能：} 用 ONNX Runtime 测试模型在实际部署时的推理速度。

\textbf{\code{main()} 流程：}
\begin{enumerate}
    \item 验证 ONNX：\code{onnx.load()} + \code{onnx.checker.check\_model()}
    \item 创建会话：\code{ort.InferenceSession()}（CPUExecutionProvider）
    \item 预热：10 次推理（排除初始化开销）
    \item 基准：100 次推理，\code{time.perf\_counter()} 高精度计时
    \item 统计：均值、标准差、P95、吞吐量
    \item 对比：若 INT8 ONNX 存在则自动 FP32 vs INT8 对比
\end{enumerate}

% =================================================================
\section{从零到一复现}

\begin{enumerate}
    \item \textbf{安装依赖：}
    \begin{lstlisting}[language=bash]
pip install torch torchvision numpy scipy scikit-learn onnx onnxruntime
    \end{lstlisting}

    \item \textbf{训练模型（自动下载数据集）：}
    \begin{lstlisting}[language=bash]
python train.py --epochs 30
    \end{lstlisting}

    \item \textbf{量化并导出 ONNX：}
    \begin{lstlisting}[language=bash]
python quantize_export.py
    \end{lstlisting}

    \item \textbf{推理性能基准：}
    \begin{lstlisting}[language=bash]
python inference_benchmark.py --onnx checkpoints/model_fp32.onnx --iters 100
    \end{lstlisting}
\end{enumerate}

% =================================================================
\section{数据流总览}

\begin{verbatim}
UCI HAR Dataset (.txt 文件)
        │
        ▼
data_loader.py ──── STFT 谱图 (3,224,224) ──→ VisionEncoder ──→ 576维
        │                                                           │
        └────── 原始IMU信号 (6,128) ──→ IMUEncoder ──→ 128维     │
                                                                    ▼
                                                     torch.cat → 704维
                                                                    │
                                                                    ▼
                                                      Fusion MLP (704→256→6)
                                                                    │
                                                                    ▼
                                                              6类活动输出
                                                                    │
                                               ┌────────────────────┘
                                               ▼
                               train.py ──→ best_model.pth (87.61%)
                                               │
                                               ▼
                               quantize_export.py ──→ model_fp32.onnx
                                                    ──→ model_int8.pth
                                               │
                                               ▼
                               inference_benchmark.py ──→ 4.83ms/次
\end{verbatim}

% =================================================================
\newpage
\section*{附录 A：依赖清单}

\begin{lstlisting}
torch>=2.0.0
torchvision>=0.15.0
numpy>=1.24.0
scipy>=1.10.0
scikit-learn>=1.3.0
pandas>=2.0.0
matplotlib>=3.7.0
onnx>=1.14.0
onnxruntime>=1.15.0
\end{lstlisting}

\section*{附录 B：6 类活动标签}

\begin{table}[h]
\centering
\begin{tabular}{@{}cl@{}}
\toprule
\textbf{索引} & \textbf{活动} \\
\midrule
0 & walking \\
1 & walking\_upstairs \\
2 & walking\_downstairs \\
3 & sitting \\
4 & standing \\
5 & laying \\
\bottomrule
\end{tabular}
\end{table}

\section*{附录 C：6 通道传感器信号}

\begin{table}[h]
\centering
\begin{tabular}{@{}cll@{}}
\toprule
\textbf{索引} & \textbf{信号名} & \textbf{传感器} \\
\midrule
0 & body\_acc\_x  & 加速度计 X 轴 \\
1 & body\_acc\_y  & 加速度计 Y 轴 \\
2 & body\_acc\_z  & 加速度计 Z 轴 \\
3 & body\_gyro\_x & 陀螺仪 X 轴 \\
4 & body\_gyro\_y & 陀螺仪 Y 轴 \\
5 & body\_gyro\_z & 陀螺仪 Z 轴 \\
\bottomrule
\end{tabular}
\end{table}

\end{document}
